

\hypertarget{fixture-chicago17-ad-parity}{%
\section{Fixture: chicago17-ad
parity}\label{fixture-chicago17-ad-parity}}

The foundational treatment of pragmatic markers in contact situations
came from Smith (2010), whose framework has shaped subsequent work.
Recent quantitative evidence reinforced the general claim (Smith 2010,
22). Classroom-discourse data produced complementary findings in Jones
and Cruz (2015), while the urban multilingual synthesis (Able, Baker,
and Chen 2019) broadened the empirical base.

A consolidated view across the three main lines of evidence is now
available (Smith 2010; Jones and Cruz 2015; Able, Baker, and Chen 2019).

Institutional reporting on classroom multilingualism adds context from
outside the academic literature (Modern Language Association 2021).
Undated industry commentary touches on similar themes (Perennial, n.d.).

\hypertarget{bibliography}{%
\section*{References}\label{bibliography}}
\addcontentsline{toc}{section}{References}

\hypertarget{refs}{}
\begin{CSLReferences}{1}{0}
\leavevmode\vadjust pre{\hypertarget{ref-able-baker-chen-2019}{}}%
Able, Amir, Bea Baker, and Chao Chen. 2019. {``Multilingual Practices in
Urban Contexts.''} \emph{Language in Society} 48 (1): 55--80.

\leavevmode\vadjust pre{\hypertarget{ref-jones-cruz-2015}{}}%
Jones, Robert, and Elena Cruz. 2015. {``Code-Switching in Classroom
Discourse.''} \emph{Applied Linguistics} 36 (2): 84--102.

\leavevmode\vadjust pre{\hypertarget{ref-mla-2021}{}}%
Modern Language Association. 2021. {``How Do I Cite a Source with
Multiple Authors?''} MLA Style Center.
\url{https://style.mla.org/citing-multiple-authors/}.

\leavevmode\vadjust pre{\hypertarget{ref-perennial-nd}{}}%
Perennial, Paul. n.d. {``Undated Commentary.''} Perennial Blog.
\url{https://example.com/undated}.

\leavevmode\vadjust pre{\hypertarget{ref-smith-2010}{}}%
Smith, Jane A. 2010. {``Pragmatic Markers in Contact Situations.''}
\emph{Journal of Pragmatics} 42 (3): 12--34.
\url{https://doi.org/10.1016/j.pragma.2010.01.005}.

\end{CSLReferences}

